% cap03/3.4_poblacion_muestra.tex
\section{Poblaci\'on y Muestra}

Para garantizar un escrutinio t\'ecnico riguroso y una comprobaci\'on arquitect\'onica inquebrantable, el presente trabajo ingenieril adopta un enfoque de muestreo puramente algor\'itmico y de software. Se descarta el tradicional muestreo perceptual humano, enfocando el dise\~no muestral en los componentes de m\'aquina y la arquitectura del sistema.

\subsection{Poblaci\'on y Muestra T\'ecnica}

La poblaci\'on t\'ecnica est\'a constituida por el universo absoluto de m\'odulos de \textit{Backend}, los \textit{endpoints} de la API RESTful, las sentencias de ruteo en \textit{PostgreSQL}, y los vectores de propagaci\'on de \textit{WebSockets}. 

Dado el rigor exigido por la ingenier\'ia de software de alta disponibilidad, la muestra t\'ecnica se define como un censo absoluto. Es decir, el $100\%$ de los m\'odulos (poblaci\'on = muestra) son sometidos a las auditor\'ias de \textit{White-Box Testing} y an\'alisis de cobertura de c\'odigo. No se puede tolerar matem\'aticamente que un s\'olo \textit{endpoint} quede fuera de la validaci\'on transaccional y unitaria, ya que un \'unico fallo en la aplicaci\'on de \textit{Row-Level Security} implicar\'ia el colapso de la seguridad del sistema y una posible exfiltraci\'on de datos.
